\section{Data collection}

\textbf{Questionnaires}

We collected data from 2 questionnaires - 1 taken before and 1 after filling out the courses. The answers were anonymized, and it is not possible to trace an answer to a specific user. We only knew which course variant the answer belongs to.

Unless stated otherwise, the questions had fixed answer selection on a scale from 0-10, where 0 means strongly disagree and 10 means strongly agree.

The first questionnaire contained questions trying to assess the students' baseline attitudes and self-reported skill level:	

\begin{itemize}
    \item How would you rate your UML skills? (Likert 0-10)
    \item What is your prior experience with UML? (Multiple choice)
    \item How interested are you in software modeling? (Likert 0-10)
    \item What is your attitude towards playful study materials? (Likert 0-10)
    \item What is your attitude towards competitive activities? (Likert 0-10)
\end{itemize}

The second questionnaire had significantly more questions, but they were still aimed at the students' attitudes towards e-learning, their satisfaction with the experience and self-reported skill level, including a number of open-ended questions. Full questionnaire is provided in Appendix~\ref{appendix:questionnaire}.\todo{ensure this is correct}

\begin{itemize}
    \item How would you rate your UML skills? (Likert 0-10)
    \item I would like to encounter more e-learning materials during my studies at \ac{fiit}. (Likert 0-10)
    \item This course helped me develop a better understanding of the UML class diagram. (Likert 0-10)
    \item I would voluntarily take a similar course. (Likert 0-10)
    \item The course kept me engaged. (Likert 0-10)
\end{itemize}

\textbf{Reports}

We collected data regarding the students' objective performance. Names and email addresses were omitted from the reports, but we did keep the user identifiers. There are 2 variants of each mentioned report - one for the competitive course version and one for the playful course. That means though no feature describes which course a data point belongs to, the data are separated by course implicitly. The only exception is Average\_score\_per\_quiz\_per\_course, which contain courseid as a feature.

Each question is categorized using tags, allowing us to examine the scores across various topics. We have average scores for each exercise and the overall number of attempts at them. This data was aggregated across all users divided by course variant.

The introductory and final assessment quizzes contain very similar exercises. That makes the difference between their relative scores a reasonable measure of progress of individual students. This data were collected per individual.

Similarly, we examined the performance in the first half of the courses' quizzes compared to the second half, and the students' final grades.

\begin{itemize}
    \item Most difficult question (describes individual exercises)
        \begin{itemize}
            \item question\_name - name of the exercise
            \item quiz\_name - name of the quiz in which the exercise is located
            \item tags - categories of the exercise
            \item attempts - total number of attempts at the exercise across all users in the course variant
            \item avg\_score\_percent - average score in the exercise normalized to a scale of 0-1
            \item variability - standard deviation of the score in the exercise
            \item difficulty\_label - derived feature, describes an exercise's difficulty based on avg\_score\_percent
        \end{itemize}
    \item Average\_score\_per\_quiz\_per\_course
        \begin{itemize}
            \item course\_id - identifier of the course variant
            \item quiz\_name - name of the quiz
            \item avg\_score\_percent - average score in the quiz across all users in the course variant
        \end{itemize}
    \item Final grade across all users
        \begin{itemize}
            \item finalgrade - the final grade, calculated as a sum of points gained across quizzes in the course. Quizzes had different weights - the final assessment quiz contributed the most significantly
        \end{itemize}
    \item First half VS second half
        \begin{itemize}
            \item first\_group\_score - score in the first 3 quizzes (expressed in percents)
            \item second\_group\_score - score in the last 2 quizzes (expressed in percents)
            \item improvement - calculated as second\_group\_score - first\_group\_score
        \end{itemize}
    \item Intro vs final quiz improvement
        \begin{itemize}
            \item intro\_score - normalized score in the Introductory assessment quiz on scale 0-1
            \item final\_score - normalized score in the Final assessment quiz on scale 0-1
            \item improvement - calculated as final\_score - intro\_score
        \end{itemize}
\end{itemize}

Full reports are provided in Appendix~\ref{appendix:reports}.\todo{ensure this is correct and anonymized correctly}

% The data was collected per-student in most cases, though in some cases we only used per-quiz aggregations.

\textbf{Moodle logs}

Moodle stores data about user behavior, such as a user's activity start and end times, interactions with the system, or number of attempts for activities. We gain insights into how many times they viewed theory sections and how many times they interacted with quizes (most of those interactions are reattempts at exercises). We collected this data for each user separately.

Measuring time taken to go through the course is unreliable, since the material is self-paced and students could take breaks in between course activities. For this reason, we approximated the time spent by only including quiz attempts. The data are still noisy, because the students could still take breaks in the middle of a quiz. But that is less likely and left us with only a handful of outliers.

The most important features of the resulting datasets were:
\begin{itemize}
    \item Interactions 		
        \begin{itemize}
            \item interaction\_type - name of the interaction type (theory/quiz/feedback)
            \item activity\_name - name of the activity (Introductory assessment quiz, Classes and objects theory page...)
            \item interactions - number of interactions with an activity
        \end{itemize}
    \item activity\_time\_sums
        \begin{itemize}
            \item total\_time\_seconds - total time spent on quizzes in seconds
            \item total\_attempts - number of attempts at quizzes in total (5 at minimum)
        \end{itemize}
\end{itemize}

Full reports are provided in Appendix~\ref{appendix:moodle_logs}.\todo{ensure this is correct}